\editionPage[\pdfbookmark[0]{Défis d’approfondissement}{bm-defis-debut-parent}]{defis-debut}{Vingt-et-un défis}{Un défi par chapitre · Parcours facultatif}{Chercher avant de consulter les deux indications.}
\begin{blocEdition}{defis-debut}
Ces problèmes demandent davantage d’initiative. Ils utilisent les outils du livre ; les prolongements sont définis dans l’énoncé. Les défis sont facultatifs et approfondissent les notions du programme.\par\bigskip
\begin{savoirfaire}[Une recherche qui avance]Faire un dessin, tester un petit cas, organiser des essais dans un tableau, garder ses traces. Ouvrir la première indication seulement après une vraie recherche ; revenir au problème avant de lire la seconde.\end{savoirfaire}
\begin{savoirfaire}[Après la solution]Repérer l’idée décisive et la refaire sur une feuille vierge, puis raconter la solution à quelqu’un. Expliquer pourquoi cela fonctionne compte autant que trouver.\end{savoirfaire}
\end{blocEdition}
\editionPage{defi-ex-1}{D01 · Défi}{Numéroter tout un livre}{Approfondissement · 20 à 45 minutes indicatives}
\begin{blocEdition}{defi-ex-1}
Un livre compte 150 pages, numérotées de 1 à 150. Combien de chiffres l’imprimeur écrit-il en tout pour numéroter le livre ?\par\vspace{5mm}{\otLight\fontsize{8}{10}\selectfont\color{otMuted}Espace de recherche}\par\par\vspace{7mm}{\color{otRule!65}\hrule height .25pt}\par\vspace{7mm}{\color{otRule!65}\hrule height .25pt}\par\vspace{7mm}{\color{otRule!65}\hrule height .25pt}\par\vspace{7mm}{\color{otRule!65}\hrule height .25pt}\par\vspace{7mm}{\color{otRule!65}\hrule height .25pt}\par\vspace{7mm}{\color{otRule!65}\hrule height .25pt}
\end{blocEdition}
\editionPage{defi-ex-2}{D02 · Défi}{Le carré magique}{Approfondissement · 20 à 45 minutes indicatives}
\begin{blocEdition}{defi-ex-2}
Compléter un carré de neuf cases avec les nombres 1 à 9, chacun une seule fois, pour que chaque ligne, chaque colonne et chaque diagonale totalisent 15. La case centrale vaut 5 et la case en haut à gauche vaut 2.\par\vspace{5mm}{\otLight\fontsize{8}{10}\selectfont\color{otMuted}Espace de recherche}\par\par\vspace{7mm}{\color{otRule!65}\hrule height .25pt}\par\vspace{7mm}{\color{otRule!65}\hrule height .25pt}\par\vspace{7mm}{\color{otRule!65}\hrule height .25pt}\par\vspace{7mm}{\color{otRule!65}\hrule height .25pt}\par\vspace{7mm}{\color{otRule!65}\hrule height .25pt}\par\vspace{7mm}{\color{otRule!65}\hrule height .25pt}
\end{blocEdition}
\editionPage{defi-ex-3}{D03 · Défi}{Plier, plier, plier}{Approfondissement · 20 à 45 minutes indicatives}
\begin{blocEdition}{defi-ex-3}
Une grande affiche est pliée en deux, puis encore en deux, et ainsi de suite. Après un pli, elle a 2 épaisseurs. Combien d’épaisseurs après 5 plis ? Combien de plis pour dépasser 100 épaisseurs ?\par\vspace{5mm}{\otLight\fontsize{8}{10}\selectfont\color{otMuted}Espace de recherche}\par\par\vspace{7mm}{\color{otRule!65}\hrule height .25pt}\par\vspace{7mm}{\color{otRule!65}\hrule height .25pt}\par\vspace{7mm}{\color{otRule!65}\hrule height .25pt}\par\vspace{7mm}{\color{otRule!65}\hrule height .25pt}\par\vspace{7mm}{\color{otRule!65}\hrule height .25pt}\par\vspace{7mm}{\color{otRule!65}\hrule height .25pt}
\end{blocEdition}
\editionPage{defi-ex-4}{D04 · Défi}{Cent œufs, douze boîtes}{Approfondissement · 20 à 45 minutes indicatives}
\begin{blocEdition}{defi-ex-4}
Une fermière range exactement 100 œufs dans 12 boîtes pleines. Elle n’a que des boîtes de 6 œufs et des boîtes de 10 œufs. Combien de boîtes de chaque sorte utilise-t-elle ?\par\vspace{5mm}{\otLight\fontsize{8}{10}\selectfont\color{otMuted}Espace de recherche}\par\par\vspace{7mm}{\color{otRule!65}\hrule height .25pt}\par\vspace{7mm}{\color{otRule!65}\hrule height .25pt}\par\vspace{7mm}{\color{otRule!65}\hrule height .25pt}\par\vspace{7mm}{\color{otRule!65}\hrule height .25pt}\par\vspace{7mm}{\color{otRule!65}\hrule height .25pt}\par\vspace{7mm}{\color{otRule!65}\hrule height .25pt}
\end{blocEdition}
\editionPage{defi-ex-5}{D05 · Défi}{La tour triangulaire}{Approfondissement · 20 à 45 minutes indicatives}
\begin{blocEdition}{defi-ex-5}
On empile des cubes en escalier : 1 cube à l’étape 1, puis 1+2=3 cubes à l’étape 2, puis 1+2+3=6 cubes à l’étape 3. Combien de cubes à l’étape 8 ? À quelle étape utilise-t-on exactement 45 cubes ?\par\vspace{5mm}{\otLight\fontsize{8}{10}\selectfont\color{otMuted}Espace de recherche}\par\par\vspace{7mm}{\color{otRule!65}\hrule height .25pt}\par\vspace{7mm}{\color{otRule!65}\hrule height .25pt}\par\vspace{7mm}{\color{otRule!65}\hrule height .25pt}\par\vspace{7mm}{\color{otRule!65}\hrule height .25pt}\par\vspace{7mm}{\color{otRule!65}\hrule height .25pt}\par\vspace{7mm}{\color{otRule!65}\hrule height .25pt}
\end{blocEdition}
\editionPage{defi-ex-6}{D06 · Défi}{Le drapeau étoilé}{Approfondissement · 20 à 45 minutes indicatives}
\begin{blocEdition}{defi-ex-6}
Un drapeau rectangulaire est partagé en trois bandes verticales de même largeur. La moitié de la bande du milieu est couverte d’étoiles. Quelle fraction du drapeau les étoiles couvrent-elles ?\par\vspace{5mm}{\otLight\fontsize{8}{10}\selectfont\color{otMuted}Espace de recherche}\par\par\vspace{7mm}{\color{otRule!65}\hrule height .25pt}\par\vspace{7mm}{\color{otRule!65}\hrule height .25pt}\par\vspace{7mm}{\color{otRule!65}\hrule height .25pt}\par\vspace{7mm}{\color{otRule!65}\hrule height .25pt}\par\vspace{7mm}{\color{otRule!65}\hrule height .25pt}\par\vspace{7mm}{\color{otRule!65}\hrule height .25pt}
\end{blocEdition}
\editionPage{defi-ex-7}{D07 · Défi}{La course des fourmis}{Approfondissement · 20 à 45 minutes indicatives}
\begin{blocEdition}{defi-ex-7}
Trois fourmis avancent sur une même règle, depuis le zéro. Fourmi A est à $\frac54$ d’unité, fourmi B à $\frac64$ et fourmi C à $\frac78$. Qui mène la course ? Qui a dépassé l’unité ?\par\vspace{5mm}{\otLight\fontsize{8}{10}\selectfont\color{otMuted}Espace de recherche}\par\par\vspace{7mm}{\color{otRule!65}\hrule height .25pt}\par\vspace{7mm}{\color{otRule!65}\hrule height .25pt}\par\vspace{7mm}{\color{otRule!65}\hrule height .25pt}\par\vspace{7mm}{\color{otRule!65}\hrule height .25pt}\par\vspace{7mm}{\color{otRule!65}\hrule height .25pt}\par\vspace{7mm}{\color{otRule!65}\hrule height .25pt}
\end{blocEdition}
\editionPage{defi-ex-8}{D08 · Défi}{Le nombre au code secret}{Approfondissement · 20 à 45 minutes indicatives}
\begin{blocEdition}{defi-ex-8}
Je suis un nombre décimal qui s’écrit avec un chiffre des unités et un chiffre des dixièmes seulement. Mon chiffre des dixièmes est le double de mon chiffre des unités, et la somme des deux chiffres vaut 9. Qui suis-je ?\par\vspace{5mm}{\otLight\fontsize{8}{10}\selectfont\color{otMuted}Espace de recherche}\par\par\vspace{7mm}{\color{otRule!65}\hrule height .25pt}\par\vspace{7mm}{\color{otRule!65}\hrule height .25pt}\par\vspace{7mm}{\color{otRule!65}\hrule height .25pt}\par\vspace{7mm}{\color{otRule!65}\hrule height .25pt}\par\vspace{7mm}{\color{otRule!65}\hrule height .25pt}\par\vspace{7mm}{\color{otRule!65}\hrule height .25pt}
\end{blocEdition}
\editionPage{defi-ex-9}{D09 · Défi}{Entre deux graduations}{Approfondissement · 20 à 45 minutes indicatives}
\begin{blocEdition}{defi-ex-9}
Trouver tous les nombres à deux chiffres après la virgule strictement compris entre $7{,}2$ et $7{,}3$. Combien y en a-t-il ? Et entre $7{,}25$ et $7{,}3$ ?\par\vspace{5mm}{\otLight\fontsize{8}{10}\selectfont\color{otMuted}Espace de recherche}\par\par\vspace{7mm}{\color{otRule!65}\hrule height .25pt}\par\vspace{7mm}{\color{otRule!65}\hrule height .25pt}\par\vspace{7mm}{\color{otRule!65}\hrule height .25pt}\par\vspace{7mm}{\color{otRule!65}\hrule height .25pt}\par\vspace{7mm}{\color{otRule!65}\hrule height .25pt}\par\vspace{7mm}{\color{otRule!65}\hrule height .25pt}
\end{blocEdition}
\editionPage{defi-ex-10}{D10 · Défi}{Cinq euros, huit pièces}{Approfondissement · 20 à 45 minutes indicatives}
\begin{blocEdition}{defi-ex-10}
Payer exactement 5 euros avec exactement 8 pièces, en n’utilisant que des pièces de 2 euros, de 1 euro et de 50 centimes. Trouver la solution et expliquer pourquoi c’est la seule.\par\vspace{5mm}{\otLight\fontsize{8}{10}\selectfont\color{otMuted}Espace de recherche}\par\par\vspace{7mm}{\color{otRule!65}\hrule height .25pt}\par\vspace{7mm}{\color{otRule!65}\hrule height .25pt}\par\vspace{7mm}{\color{otRule!65}\hrule height .25pt}\par\vspace{7mm}{\color{otRule!65}\hrule height .25pt}\par\vspace{7mm}{\color{otRule!65}\hrule height .25pt}\par\vspace{7mm}{\color{otRule!65}\hrule height .25pt}
\end{blocEdition}
\editionPage{defi-ex-11}{D11 · Défi}{La corde à border}{Approfondissement · 20 à 45 minutes indicatives}
\begin{blocEdition}{defi-ex-11}
Avec une corde de 24 m, on forme exactement le tour d’une figure. Quel est le côté du carré que l’on peut former ? Et si l’on forme un rectangle de 4 m de large, quelle est sa longueur ?\par\vspace{5mm}{\otLight\fontsize{8}{10}\selectfont\color{otMuted}Espace de recherche}\par\par\vspace{7mm}{\color{otRule!65}\hrule height .25pt}\par\vspace{7mm}{\color{otRule!65}\hrule height .25pt}\par\vspace{7mm}{\color{otRule!65}\hrule height .25pt}\par\vspace{7mm}{\color{otRule!65}\hrule height .25pt}\par\vspace{7mm}{\color{otRule!65}\hrule height .25pt}\par\vspace{7mm}{\color{otRule!65}\hrule height .25pt}
\end{blocEdition}
\editionPage{defi-ex-12}{D12 · Défi}{Quatre litres exactement}{Approfondissement · 20 à 45 minutes indicatives}
\begin{blocEdition}{defi-ex-12}
Au bord d’une fontaine, on dispose d’un bidon de 5 L et d’un bidon de 3 L, sans graduation. Comment obtenir exactement 4 L d’eau dans le grand bidon ?\par\vspace{5mm}{\otLight\fontsize{8}{10}\selectfont\color{otMuted}Espace de recherche}\par\par\vspace{7mm}{\color{otRule!65}\hrule height .25pt}\par\vspace{7mm}{\color{otRule!65}\hrule height .25pt}\par\vspace{7mm}{\color{otRule!65}\hrule height .25pt}\par\vspace{7mm}{\color{otRule!65}\hrule height .25pt}\par\vspace{7mm}{\color{otRule!65}\hrule height .25pt}\par\vspace{7mm}{\color{otRule!65}\hrule height .25pt}
\end{blocEdition}
\editionPage{defi-ex-13}{D13 · Défi}{Les pièces du carré}{Approfondissement · 20 à 45 minutes indicatives}
\begin{blocEdition}{defi-ex-13}
Un carré de 16 carreaux est découpé en trois pièces : un grand triangle qui couvre la moitié du carré, et deux pièces qui se partagent le reste à égalité. Donner l’aire de chaque pièce, en carreaux.\par\vspace{5mm}{\otLight\fontsize{8}{10}\selectfont\color{otMuted}Espace de recherche}\par\par\vspace{7mm}{\color{otRule!65}\hrule height .25pt}\par\vspace{7mm}{\color{otRule!65}\hrule height .25pt}\par\vspace{7mm}{\color{otRule!65}\hrule height .25pt}\par\vspace{7mm}{\color{otRule!65}\hrule height .25pt}\par\vspace{7mm}{\color{otRule!65}\hrule height .25pt}\par\vspace{7mm}{\color{otRule!65}\hrule height .25pt}
\end{blocEdition}
\editionPage{defi-ex-14}{D14 · Défi}{L’horloge pressée}{Approfondissement · 20 à 45 minutes indicatives}
\begin{blocEdition}{defi-ex-14}
Une horloge déréglée avance de 3 minutes toutes les heures. On la met à l’heure à 8 h du matin.\begin{questions}
\item Quelle heure affiche-t-elle quand il est réellement 14 h ?
\item Quelle avance aura-t-elle prise après 24 heures ?
\item Quand l’horloge affiche 20 h 36 le même jour, quelle heure est-il réellement ?
\end{questions}\par\vspace{5mm}{\otLight\fontsize{8}{10}\selectfont\color{otMuted}Espace de recherche}\par\par\vspace{7mm}{\color{otRule!65}\hrule height .25pt}\par\vspace{7mm}{\color{otRule!65}\hrule height .25pt}\par\vspace{7mm}{\color{otRule!65}\hrule height .25pt}\par\vspace{7mm}{\color{otRule!65}\hrule height .25pt}\par\vspace{7mm}{\color{otRule!65}\hrule height .25pt}\par\vspace{7mm}{\color{otRule!65}\hrule height .25pt}
\end{blocEdition}
\editionPage{defi-ex-15}{D15 · Défi}{Le trésor du jardin}{Approfondissement · 20 à 45 minutes indicatives}
\begin{blocEdition}{defi-ex-15}
Exécuter ce programme sur une feuille : tracer un segment [AB] de 8 cm ; placer son milieu M ; tracer la perpendiculaire à (AB) passant par M ; placer sur cette perpendiculaire le point T à 3 cm de M. Le trésor T est-il à la même distance de A et de B ? Vérifier au compas et expliquer.\par\vspace{5mm}{\otLight\fontsize{8}{10}\selectfont\color{otMuted}Espace de recherche}\par\par\vspace{7mm}{\color{otRule!65}\hrule height .25pt}\par\vspace{7mm}{\color{otRule!65}\hrule height .25pt}\par\vspace{7mm}{\color{otRule!65}\hrule height .25pt}\par\vspace{7mm}{\color{otRule!65}\hrule height .25pt}\par\vspace{7mm}{\color{otRule!65}\hrule height .25pt}\par\vspace{7mm}{\color{otRule!65}\hrule height .25pt}
\end{blocEdition}
\editionPage{defi-ex-16}{D16 · Défi}{Le pliage du carré}{Approfondissement · 20 à 45 minutes indicatives}
\begin{blocEdition}{defi-ex-16}
On plie un carré de 8 cm de côté le long d’une diagonale. Quelle figure obtient-on ? Son périmètre est-il égal à la moitié du périmètre du carré ? Expliquer sans mesurer la diagonale.\par\vspace{5mm}{\otLight\fontsize{8}{10}\selectfont\color{otMuted}Espace de recherche}\par\par\vspace{7mm}{\color{otRule!65}\hrule height .25pt}\par\vspace{7mm}{\color{otRule!65}\hrule height .25pt}\par\vspace{7mm}{\color{otRule!65}\hrule height .25pt}\par\vspace{7mm}{\color{otRule!65}\hrule height .25pt}\par\vspace{7mm}{\color{otRule!65}\hrule height .25pt}\par\vspace{7mm}{\color{otRule!65}\hrule height .25pt}
\end{blocEdition}
\editionPage{defi-ex-17}{D17 · Défi}{Les mots du miroir}{Approfondissement · 20 à 45 minutes indicatives}
\begin{blocEdition}{defi-ex-17}
Certaines lettres majuscules ont un axe de symétrie vertical : A, H, I, M, O, T, U, V, W, X, Y. Expliquer pourquoi le mot OTTO se lit encore dans un miroir vertical, puis trouver un autre mot qui résiste au miroir.\par\vspace{5mm}{\otLight\fontsize{8}{10}\selectfont\color{otMuted}Espace de recherche}\par\par\vspace{7mm}{\color{otRule!65}\hrule height .25pt}\par\vspace{7mm}{\color{otRule!65}\hrule height .25pt}\par\vspace{7mm}{\color{otRule!65}\hrule height .25pt}\par\vspace{7mm}{\color{otRule!65}\hrule height .25pt}\par\vspace{7mm}{\color{otRule!65}\hrule height .25pt}\par\vspace{7mm}{\color{otRule!65}\hrule height .25pt}
\end{blocEdition}
\editionPage{defi-ex-18}{D18 · Défi}{Le grand cube peint}{Approfondissement · 20 à 45 minutes indicatives}
\begin{blocEdition}{defi-ex-18}
Vingt-sept petits cubes sont assemblés en un grand cube de 3 sur 3 sur 3, puis toute la surface du grand cube est peinte. Combien de petits cubes ont exactement 3 faces peintes ? Exactement 2 ? Aucune ?\par\vspace{5mm}{\otLight\fontsize{8}{10}\selectfont\color{otMuted}Espace de recherche}\par\par\vspace{7mm}{\color{otRule!65}\hrule height .25pt}\par\vspace{7mm}{\color{otRule!65}\hrule height .25pt}\par\vspace{7mm}{\color{otRule!65}\hrule height .25pt}\par\vspace{7mm}{\color{otRule!65}\hrule height .25pt}\par\vspace{7mm}{\color{otRule!65}\hrule height .25pt}\par\vspace{7mm}{\color{otRule!65}\hrule height .25pt}
\end{blocEdition}
\editionPage{defi-ex-19}{D19 · Défi}{La photo agrandie}{Approfondissement · 20 à 45 minutes indicatives}
\begin{blocEdition}{defi-ex-19}
Une photo mesure 10 cm sur 15 cm. On l’agrandit pour que le petit côté mesure 30 cm. Quelle est la nouvelle longueur du grand côté ? L’agrandissement tiendra-t-il dans un cadre de 30 cm sur 40 cm ?\par\vspace{5mm}{\otLight\fontsize{8}{10}\selectfont\color{otMuted}Espace de recherche}\par\par\vspace{7mm}{\color{otRule!65}\hrule height .25pt}\par\vspace{7mm}{\color{otRule!65}\hrule height .25pt}\par\vspace{7mm}{\color{otRule!65}\hrule height .25pt}\par\vspace{7mm}{\color{otRule!65}\hrule height .25pt}\par\vspace{7mm}{\color{otRule!65}\hrule height .25pt}\par\vspace{7mm}{\color{otRule!65}\hrule height .25pt}
\end{blocEdition}
\editionPage{defi-ex-20}{D20 · Défi}{Le tableau effacé}{Approfondissement · 20 à 45 minutes indicatives}
\begin{blocEdition}{defi-ex-20}
Le tableau des inscriptions aux deux clubs de l’école est en partie effacé. On sait : 26 élèves au total en CM1, 24 élèves au club théâtre en tout, 55 élèves en tout, et 14 élèves de CM2 au club jeux. Reconstituer tout le tableau : lignes CM1 et CM2, colonnes théâtre et jeux.\par\vspace{5mm}{\otLight\fontsize{8}{10}\selectfont\color{otMuted}Espace de recherche}\par\par\vspace{7mm}{\color{otRule!65}\hrule height .25pt}\par\vspace{7mm}{\color{otRule!65}\hrule height .25pt}\par\vspace{7mm}{\color{otRule!65}\hrule height .25pt}\par\vspace{7mm}{\color{otRule!65}\hrule height .25pt}\par\vspace{7mm}{\color{otRule!65}\hrule height .25pt}\par\vspace{7mm}{\color{otRule!65}\hrule height .25pt}
\end{blocEdition}
\editionPage{defi-ex-21}{D21 · Défi}{Le sac mystère}{Approfondissement · 20 à 45 minutes indicatives}
\begin{blocEdition}{defi-ex-21}
Un sac opaque contient 10 boules, rouges ou bleues. Tirer une boule rouge est « probable » (plus d’une chance sur deux), et tirer une bleue est possible, mais peu probable. Combien de boules rouges le sac peut-il contenir ? Donner toutes les possibilités.\par\vspace{5mm}{\otLight\fontsize{8}{10}\selectfont\color{otMuted}Espace de recherche}\par\par\vspace{7mm}{\color{otRule!65}\hrule height .25pt}\par\vspace{7mm}{\color{otRule!65}\hrule height .25pt}\par\vspace{7mm}{\color{otRule!65}\hrule height .25pt}\par\vspace{7mm}{\color{otRule!65}\hrule height .25pt}\par\vspace{7mm}{\color{otRule!65}\hrule height .25pt}\par\vspace{7mm}{\color{otRule!65}\hrule height .25pt}
\end{blocEdition}
